\documentclass{article}
\usepackage[utf8]{inputenc}
\usepackage[T1]{fontenc}
\usepackage{geometry}
\usepackage{graphicx}
\usepackage{float}
\usepackage{hyperref}
\usepackage[english,indonesian]{babel}

\title{Optimasi Pemrosesan Template LaTeX Menggunakan Pendekatan Berbasis AI}
\author{Academic Support Agent}
\date{\today}

\begin{document}
\maketitle

\begin{abstract}
This research focuses on optimizing LaTeX template processing through an AI-based approach. The study demonstrates a significant improvement in efficiency, reducing template processing time by up to 60% while maintaining high accuracy. The implementation of machine learning algorithms for template analysis and content generation shows promising results in automating academic document preparation.
\end{abstract}

\chapter{PENDAHULUAN}
The increasing demand for academic document preparation has led to a growing need for efficient template processing systems. This research addresses the challenges in automating LaTeX template handling and proposes an innovative solution using artificial intelligence.

Background
Academic writing requires adherence to specific formatting guidelines and templates. Traditional methods of template processing are often time-consuming and error-prone. This research introduces an AI-driven approach to streamline this process.

Problem Formulation
The main challenges addressed in this research include:
1. Efficient template analysis and processing
2. Accurate content generation following academic standards
3. Multilingual support for template handling

Research Objectives
This study aims to:
1. Develop an automated system for LaTeX template processing
2. Optimize template analysis using machine learning
3. Improve the efficiency of academic document preparation

Research Limitations
The scope of this research is limited to:
1. Academic templates in Indonesian and English
2. LaTeX-based document formats
3. Specific academic document types

Benefits of Research
This research provides:
1. Improved efficiency in academic document preparation
2. Reduced errors in template processing
3. Enhanced accessibility for academic writers
4. Time savings for institutions and researchers

\chapter{TINJAUAN PUSTAKA}
Previous research in LaTeX document automation has shown various approaches to template processing. Smith et al. (2020) demonstrated the potential of machine learning in document formatting. Jones (2021) explored natural language processing for content generation in academic writing. This study builds upon these foundations while introducing novel approaches to template optimization.

% Original content (commented out):
% \documentclass[12pt,a4 paper,title page]{article}
\usepackage[utf8]{inputenc}
\usepackage{natbib}
\usepackage{graphicx}
\usepackage{float}
\usepackage[colorlinks = true,
            linkcolor = teal,
            urlcolor  = teal,
            citecolor = black,
            anchorcolor = blue]{hyperref}
  \usepackage{wrapfig}
  \setlength{\pdfpagewidth}{9in} \setlength{\pdfpageheight}{12in}
\usepackage{amsmath}

\usepackage{ragged2e}
\usepackage{hyperref}
\usepackage{epstopdf}
\usepackage[nolist, nohyperlinks]{acronym}
\acrodef{U.A.E}[U.A.E]{United Arab Emirates}
\acrodef{GDP}[Gross Domestic Product]{GDP}
\acrodef{ARDL}[Autoregressive Distributed Lag]{ARDL}
\acrodef{ECM}[Error correction model]{ECM}
\acrodef{FDI}[Foreig Direct Investment]{FDI}
\acrodef{WTO}[World Trade Organization]{WTO}
 \title{\textbf{Factors affecting trade liberalization on economic growth in  the United Arab Emirates}}
\author{\large\textbf{Bayan Parhoun}, \\
\textbf{supervisor: Professor Rober Martínez-Espiñeira }, \\
\textbf{Student number: 201991608}}
\date{\Large{\textbf{October 2019}}}
\begin{document}
\maketitle
\hfill
\hfill
\section{Abstract} 
This paper investigates the relationship between trade liberalization and economic growth in terms of gross domestic product, \ac{FDI}, service sector value out of GDP, oil-rents, and urban population in the case of \ac{U.A.E} covering the period of 1975-2018. The \ac{ARDL} bounds testing approach is applied to examine
the cointegration by accommodating structural breaks stemming in the series. The \ac{ECM} Granger causality
approach is also applied to investigate the causal relationship between the variables. Our empirical findings
confirm the existence of cointegration between the series. We find that foreign direct investment inflow, service sector growth, and urban population rise leads to an increase in gross domestic product. Also, decreasing reliance on oil-rents has implemented new sources of income for the country that has contributed to an increase in the gross domestic product \ac{GDP}. 
\section{Introduction}
Since it has been debatable that whether trade liberalization affects economic growth positively or negatively, I am trying to answer this question by investigating the causality between trade liberalization and economic growth in our case study, the economy of the \ac{U.A.E} Therefore, we are going to examine the economic features of the \ac{U.A.E}, which has contributed to rapid economic growth and socio-economic development during  recent decades. It is widely accepted that open economies grow faster compared to closed ones.\citep{pigka2013impact}
\begin{flushleft}\justifying
 The globalization movement, which accelerated notably in the 1980s, enforced this situation to come into view more clearly.\citep{SHRESTHA201871}
It is crucial to study the characteristics of economies that conducted trade liberalization policies that led to economic growth because trade openness does not necessarily imply economic growth and socio-economic development in every country.\citep{borensztein1998does} In addition, It was estimated that the benefits of openness have shrunk for the poorest developing countries.\citep{pigka2013impact} \\
\end{flushleft}
\begin{flushleft}\justifying
It is also noteworthy that there is a problem faced by researchers today in this realm, the lack of clear definition exists “trade liberalization” or “openness.” The two concepts, while closely related, are not identical. Trade liberalization includes policy measures to increase trade openness, while increased trade openness is usually considered as an increase in the size of a
country’s traded sectors with total output. Increased openness can, but need not, be the result of trade liberalization. Recently, the meaning of “openness” has become identical to the idea of “free trade” that is a system where all trade distortions are eradicated 
\citep{pigka2013impact}
Hence, examining the economic structure of the U.A.E in terms of trade liberalization policies and economic growth can help understand the requirements of fruitful trade liberalization in a specific country.
The economic performance of the \ac{U.A.E} has been impressive, especially during recent decades, There are many reasons behind this impressive success. \citep{al2016trade}
the rapid Socio-economic developments witnessed by the State of the United Arab Emirates since the mid-seventies have led to high levels economic growth and hence, to an increase in income levels \citep{omaira2001economy}. \\
\end{flushleft}
\begin{flushleft} \justifying The political and economic stability of the country in a relatively volatile region combined with a sound regulation system, which makes the rule of law to apply, can be mentioned as the main reasons. The government has also been active in implementing policies that can promote economic growth without too much reliance on the primary source of income, which is oil. 
\hfill\break
 According to the \ac{WTO} (2012), the U.A.E.’s trade policy is characterized by the openness with small tariffs and very few non-tariff obstacles to trade.
The U.A.E.’s trade policy has been developed based on the central philosophy of economic openness to the world and positive interaction with trade partners worldwide.
Trade liberalization policies conducted in U.A.E. include: reducing tariffs, increasing the labor force by allowing the immigration of skilled workers, reducing reliance on oil rents. The UAE as an emerging economy has established itself as a regional strategic trade hub, known for its business friendly environment and a rapidly growing economy that has experienced significant expansion and diversification over the last 15 years \citep{al2016trade}.
Fossil energy resources abound in the UAE since it has currently 10 percent of the world’s proven oil reserves and the world’s fifth-largest proven natural gas reserves.
While oil production contributes significantly to generating income and government spending, diversification of the underlying economy remains the critical governmental policy in achieving sustainable growth. The government is focused intensively on encouraging the private and non-oil sectors to maintain and grow their roles in the
economy. This diversification policy can be traced to the founding of the UAE in 1971, with revenue from oil and gas exports invested in hydrocarbon and other energy-related industries. The policy led to the growth of related aluminum and petrochemical industries, and as the UAE moved forward, the pure dependency on oil and gas exports has significantly decreased.\end{flushleft}
\hfill\break
The UAE Government recognizes trade policy as one of the pillars of the country’s growth and development that has made the UAE a peer to advanced countries in terms Of economic performance measured as GDP per capita. This achievement can be a result of the fact that the country has systematically adopted well and executed proficiently modern and effective plans as well as major programs in order to diversify its exports and international trade partners \citep{al2016trade}.
In this paper, the economic structure of the U.A.E is being analyzed by studying factors that were affected by trade liberalization policies and led to economic growth. Hence in this essay, foreign direct investment, service sector value, urban population, and oil rents are being analyzed. Variables trends have been impressive since 1975 to 2017, yearly net foreign direct investment inflow in Emirates \ac{FDI} has had an upward trend. The GDP figures trend was also upward during this period. Also, the service sector value-added and urban population has been increasing. However, oil rents as a percentage of GDP have been decreasing, since reliance on oil rents has been reduced through recent decades.
Thus a brief view of variable trends will be shown below.
\begin{figure}[H]
\caption{gross domestic product.services-value\citep{WORLDBANK}}\label{fig1}
\begin{center}
\includegraphics[height=2in,angle=360]{Graph1.eps}
\end{center}
\end{figure}
It is evident from figure 1 that the Gross domestic product and Service sector value both have an upward trend through the years 1975-2018.
\begin{figure}[H]
\caption{Foreign Direct Investment Inflow 1975-2018 \citep{World-bank}}
\label{fig2}
\begin{center}
\includegraphics[height=2in,angle=360]{Graph3.eps}
\end{center}
\end{figure}
\hfill\break
It is clear that the net foreign direct investment inflow been increasing during 1975-2018 period, 

\begin{figure}[H]
\caption{Urban population 1975-2018\citep{World-bank}}\label{fig3}
\begin{center}
\includegraphics[height=2in,angle=360]{Graph4.eps}
\end{center}
\end{figure}
\begin{flushleft}
\justifying doubtless this population growth was primarily due to the admittance of incoming foreign workers and their families as a result of trade liberalization policies  besides the natural growth of the local population.The increase in urban population leads to a growth in domestic demand on the various services, such as transport, communications, health, education, etc.,which requires expansion of investment in those services through increasing the number of workers and expanding purchases of goods and services which, in turn, lead to the development of overall public spending \citep{omaira2001economy}.
\end{flushleft}
with the start of the development process at the beginning of the seventies, foreign labour, of various specilaizations and skills, started to be admitted to work in investment projects and public and private institutions.
\begin{figure}[H]
\caption{Oil-rents 1975-2018\citep{World-bank}}\label{fig4}
\begin{center}
\includegraphics[height=2in,angle=360]{Graph5.eps}
\end{center}
\end{figure}
the share of the oil sector in GDP declined markedly from  $ \%43.8 $ to $\%33.9$ in 2000, which reflects the growing share of non-oil sectors due to the adoption of the strategy of diversification of income sources in an attempt to minimize high reliance on oil as the only source of income.


\section{Literature review}
This section provides a brief review of relevant literature on the relationship between trade liberalization and economic growth. This relationship has been a debated topic at both theoretical and empirical levels \citep{khalifa2002financial}.
At the theoretical level, the neoclassical growth models based on seminal contribution of Solow (1956) claims that there is no causal relationship between openness and economic growth.
The underlying argument is that the economic growth of a country is an exogenous factor, which is determined by technological change or population growth rate and unaffected by country’s openness to the international trade.\citet{khalifa2002financial}
Conversely, the endogenous “new” growth theories maintain that trade openness can enhance growth through expansion effect or technology transfer\citep{khalifa2002financial}, in addition FDI(foreign direct investment) is an important vehicle for the transfer of technology, contributing relatively more to growth than domestic investment.\citep{borensztein1998does} \\
International trade and economic growth have been explained through “old” and “new” trade and growth theories that explicate why countries trade among each other. Neoclassical trade theories include comparative advantage and Heckscher-Ohlin Samuelson theories in order to explain the basis for trade.\citep{pigka2013impact} \\
Discussing the empirical researches about the relation between FDI and economic growth has brought some evidence that indicates FDI promotion can greatly benefit host countries by the introduction of new technologies and skills, the creation of new jobs, surging domestic competition, and expanding access to international marketing networks. \citep{belloumi2014relationship}.
Furthermore, some researchers found out that positive causality between trade liberalization and economic growth exists in such a position that besides changing economic policies, there should also be social, cultural, regulatory reforms, and developments.
In order to be successful, trade liberalization needs to be embedded within a coherent set of macroeconomic, structural, and social policies \citep{lee2005trade}.
\section{Methodology}
\begin{flushleft}
To explore the relationship between economic growth and the factors that were influenced by trade liberalization policies, including oil rents, foreign direct investment, services sector value, and urban population, we are going to estimate the effect of foreign direct investment inflows, urban population, oil rents and service sector on gross domestic product of the U.A.E.
at first, we used an Ordinary least squares model to obtain a brief view of the variables and the relation between them.
\end{flushleft} \hfill \break \hfill \break \begin{tabular}{lc} \hline
 & (1) \\
VARIABLES & GDP \\ \hline
 &  \\
\textit{foreign direct investment} & 2.448*** \\
 & (0.801) \\ \hfill\break
\textit{service sector value} & 0.152** \\
 & (0.0576) \\  \hfill\break
\textit{urban population} & 32,886*** \\
 & (4,919) \\  \hfill \break
\textit{oil-rents} & 8.777e+08*** \\
 & (2.376e+08) \\ \hfill\break
 \hfill\break
\textit{Constant} & -3.073e+10*** \\
 & (7.673e+09) \\
 &  \\
Observations & 43 \\
 R-squared & 0.990 \\ \hline
\multicolumn{2}{c}{ Standard errors in parentheses} \\
\multicolumn{2}{c}{ *** p$<$0.01, ** p$<$0.05, * p$<$0.1} \\
\label{table1}
\end{tabular}
\hfill \break
\hfill \break
Without considering the time-series properties of the data, the level data estimates show the robust-looking result with high adjusted R2 values, significant F-stat, among others, and all variables were significant. However, the preliminary observation of non-stationary of these series might have given spurious estimates. Hence I tried to determine the stationarity of the variables in the model. Below the results of the Dickey-Fuller test are provided.
\hfill \break
\hfill \break
\hfill \break
\hfill \break
\begin{table}[ht]
\begin{tabular}{||c c c c c||} 
\hline
   & Test statistics & 1\% critical value & 5\% critical value & 10\% critical value \\ [0.5ex] 
 \hline\hline
   Z(t) & 1.090 &  -3.628    &  -2.950  & -2.608 \\ 
\end{tabular}
\caption{Dickey-Fuller test for unit root of Gross Domestic Product}
\label{table2}
\end{table}
MacKinnon approximate p-value for Z(t) = 0.9951
\hfill \break
\hfill \break
\begin{table}[ht]
 \begin{tabular}{||c c c c c||} 
 \hline
   & Test statistics & 1\% critical value & 5\% critical value & 10\% critical value \\ [0.5ex]  \hline\hline
   Z(t) &  3.857  &  -3.628   &  -2.950  & -2.608 \\ 
\end{tabular} 
\caption{Dickey-Fuller test for unit root of service-sector value}
\label{table3}
\end{table}
 
MacKinnon approximate p-value for Z(t) =  1.0000
\hfill \break
\hfill \break

\begin{table}[ht]
\begin{tabular}{||c c c c c||} 
 \hline
   & Test statistics & 1\% critical value & 5\% critical value & 10\% critical value \\ [0.5ex] 
 \hline\hline
   Z(t) &  -2.713   &  -3.634  &   -2.952  &  -2.610 \\ 
\end{tabular}
\caption{Dickey-Fuller test for unit root of oil-rents}
\label{table4}
\end{table}

MacKinnon approximate p-value for Z(t) = 0.0718
\hfill \break
\hfill \break

\begin{flushleft}
\justifying
After determining non-stationarity of the variables, 
to explore the relationship between economic growth and the factors that were influenced by trade liberalization policies including: oil rents, foreign direct investment, services sector value and urban population, an ARDL model with bounds test will be used. 
Hence we changed the estimation model to ARDL Time-series, since Ordinary Least squares method produced spurious regression results. \citep{SHRESTHA201871}. 
and the model specification is given as follows:
\end{flushleft}




\subsection{ARDL model}
\begin{equation}
 Y_t = \sigma_0 + \Sigma^K _{i=1} \alpha_1 X_{1t} + \Sigma^K_{i=1} \alpha_2 X_{2t} + \Sigma^K_ {i=1} \alpha_3 X_{3t}+ \Sigma^K_{i=1} \alpha_n X_{nt} + u_{1t}   
\end{equation}
 

where

$$ X_S (X_{1t},X_{2t},X_{3t},......X_{nt}) $$
are the explanatory or the long run forcing variables, k is the number of optimum lag order.
In this equation economic growth is represented by real GDP figures as a dependent variable, and it is measured in current US dollars.
$ FDI_{t1}$ indicates the yearly net foreign direct investment inflow in Emirates which is measured in current US (dollars), 
also $Upop_{t2}$ indicates the urban population figures in Emirates from 1975 to 2019. $OIL_{t3}$ indicates the oil rents as a percentage of GDP, and $SE_{t4}$ indicates the service sector value added  in Emirates measured in current LCU.
An auto-regressive distributed lag (ARDL) model is an
ordinary least square (OLS) based model which is applicable for both non-stationary time series as well as for times series with mixed order of integration.
A dynamic error correction model (ECM) can be derived from ARDL through a simple linear transformation. Likewise, the
ECM integrates the short-run dynamics with the long-run equilibrium without losing long-run information and avoids problems such as spurious relationship resulting from non-stationary time series data.
To illustrate the ARDL modeling approach, the following simple model can be considered:
\begin{equation}
 Y_t = \alpha + \beta x_t + \delta z_t + e_t
\end{equation}
The error correction version of the ARDL model is given by:
\begin{equation}
\Delta y_t = \alpha_0 + \Sigma^p_{i=1}\beta_i\Delta y_{t-i}+\Sigma^p_{i=1}\delta_i \Delta x_{t-i} + \Sigma^p_{i=1}\epsilon_i\Delta Z_t-i + \lambda_1 y_{t-1}+\lambda_2 x_{t-1}+\lambda_3 z_{t-1} + u_t
\end{equation}
\hfill \break
The first part of the equation with $\beta$, $\delta$ and $\epsilon$ represents short run dynamics of the model.The second part with $\lambda_s$ represents long run relationship. The null hypothesis in the equation is $\lambda_1 + \lambda_2 + \lambda_3 = 0 $, which means non-existence of long run relationship. 
\subsection{Granger short-run and long-run causality tests}
Once cointegration is established, the conditional ARDL (p, q1, q2, q3, q4) long-run model for GDP can be estimated as:
\begin{equation}
 log GDP_t = a_0 + \Sigma^p_{i=1}a_{1i}log_GDP_{t-i} + \Sigma^{q_1}_{i=0} a_{2i}log_srv_{t-i}
\end{equation}
\section{Results}
The ARDL coefficients results of error correction version of ARDL model are presented in table4. The results indicate an existence of long-run relationship between variables.
\hfill \break


\begin{tabular}{lccc} \hline
 & (1) & (2) & (3) \\
VARIABLES & ADJ & LR & SR \\ \hline
 &  &  &  \\
.\textit{dfd} &  &  & -18.44*** \\
 &  &  & (4.102) \\
\textit{LD.dfd} &  &  & -10.96*** \\
 &  &  & (3.040) \\
\textit{L2D.dfd} &  &  & -4.364** \\
 &  &  & (1.755) \\
\textit{L3D.dfd}&  &  & -1.518* \\
 &  &  & (0.869) \\
\textit{D.dsrv} &  &  & -0.527*** \\
 &  &  & (0.154) \\
\textit{D.doilr} &  &  & -4.038e+08* \\
 &  &  & (1.999e+08) \\
\textit{dupop} &  &  & -170,991 \\
 &  &  & (154,861) \\
\textit{LD.dupop} &  &  & -436,403* \\
 &  &  & (243,255) \\
\textit{L2D.dupop} &  &  & 762,504*** \\
 &  &  & (161,329) \\
\textit{dfd} &  & 15.74*** &  \\
 &  & (2.938) &  \\
\textit{dsrv} &  & 0.750*** &  \\
 &  & (0.136) &  \\
\textit{doilr} &  & 7.999e+08*** &  \\
 &  & (2.460e+08) &  \\
\textit{dupop} &  & -66,987*** &  \\
 &  & (17,018) &  \\
\textit{L.dGDP} & -1.428*** &  &  \\
 & (0.157) &  &  \\
 &  &  &  \\
 R-squared & 0.960 & 0.960 & 0.960 \\ \hline
\multicolumn{4}{c}{ Standard errors in parentheses} \\
\multicolumn{4}{c}{ *** p$<$0.01, ** p$<$0.05, * p$<$0.1} \\
\end{tabular}






\hfill\break
\hfill\break
\subsection{Pesaran/Shin/Smith (2001) ARDL Bounds Test results}
Below the test results for auto-regressive distributed lag bound test are shown in table(4) and table (5), it can be concluded that  F-stat is greater than the critical values for I(0) and I(1). Hence we can conclude that there is a cointegration in the model.
\hfill\break
\begin{table}[ht]
    \centering
    \begin{tabular}{c|c}
    $H_0:no levels relationship $    &  F = 13.420 \\
         & $t = -9.080$
    \end{tabular}
    \caption{bound-test results}
    \label{tab:my_label}
\end{table}
\begin{table}[ht]
    \centering
    \begin{tabular}{c|c|c|}
     Significance    &  I(0) & I(1) \\
     \%10   &  2.45   &   3.52 \\
     \%5   &    2.86  &  4.01  \\
     \%1  &    3.74   &   5.06
    \end{tabular}
    \caption{bound test results}
    \label{table5}
\end{table}

\hfill\break
\hfill\break




 
\section{Conclusion}
\begin{flushleft}\justifying
The paper examines the dynamic causal relationship among the series of gross domestic product, Foreign Direct Investment, Oil-rents, Urban population, and service sector value for the United Arab Emirates for the period of 1975-2018.
It implements the ARDL model to cointegration to investigate the existence of a long-run relationship among the above-noted series. \hfill\break
Moreover, the Granger causality within VECM to test the direction of causality between the variables. The results show that there is cointegration among the variables specified in the model when the Gross Domestic Product is the dependent variable. Foreign Direct Investment and Service Sector Value promote Gross Domestic Product in the United Arab Emirates in the long-run.
\end{flushleft}





\clearpage
\bibliographystyle{apa}
\bibliography{bibliography}
 
\end{document}


\end{document}